\documentclass[14pt]{article}
\usepackage[utf8]{inputenc}	% Para caracteres en español
\usepackage{amsmath,amsthm,amsfonts,amssymb,amscd}
\usepackage{multirow,booktabs}
\usepackage[table]{xcolor}
\usepackage{fullpage}
\usepackage{lastpage}
\usepackage{enumitem}
\usepackage{fancyhdr}
\usepackage{mathrsfs}
\usepackage{wrapfig}
\usepackage{setspace}
\usepackage{calc}
\usepackage{multicol}
\usepackage{cancel}
\usepackage[title]{appendix}
\usepackage[retainorgcmds]{IEEEtrantools}
\usepackage[margin=3cm]{geometry}
\usepackage{amsmath, amsthm}
\newlength{\tabcont}
\setlength{\parindent}{0.0in}
\setlength{\parskip}{0.05in}
\usepackage{empheq}
\usepackage{enumitem}  
\usepackage{framed}
\usepackage[most]{tcolorbox}
\usepackage{xcolor}
\colorlet{shadecolor}{orange!15}
\parindent 0in
\parskip 14pt
\geometry{margin=1in, headsep=0.25in}
\newtheorem{theorem}{Theorem}
\newtheorem{lemma}[theorem]{Lemma}
\newtheorem*{remark}{Remark}
\newtheorem*{solution}{Solution}
\newtheorem{corollary}[theorem]{Corollary}
\newtheorem{proposition}[theorem]{Proposition}
\theoremstyle{definition}
\newtheorem*{definition}{Definition}
\newtheorem*{exercise}{Exercise}
\newtheorem{example}[theorem]{Example}
\theoremstyle{remark}
\newtheorem{notation}{Notation}
%\AtAppendix{\counterwithin{theorem}{section}}
\setlength{\parskip}{1em}
\DeclareMathOperator{\sign}{sign}
\renewcommand{\thetheorem}{\arabic{section}.\arabic{theorem}}
\renewcommand{\theequation}{\arabic{section}.\arabic{equation}}
\usepackage[numbib,notlof,notlot,nottoc]{tocbibind}
\numberwithin{equation}{section}
\numberwithin{theorem}{section}
\usepackage{xeCJK}
    \setCJKmainfont[AutoFakeBold=3]{AR PL UKai TW}
    \XeTeXlinebreaklocale "zh"             
    \XeTeXlinebreakskip = 0pt plus 1pt


\printbibliography[heading=bibintoc,title={References}]

\begin{document}
\setcounter{section}{0}
\title{Chapter 2}

\thispagestyle{empty}
\begin{center}
{\LARGE \bf Chapter 2 in \cite{Hog10}}
\end{center}

\begin{exercise}[2.1-7]
    Let a random experiment be the casting of a pair of unbiased six-sided dice, 
    and let $X$ equal the smaller of the outcomes if they are different and the common 
    value if they are equal.
    \begin{enumerate}
    \item With reasonable assumptions, find the p.m.f. of $X$.
    \item Let $Y$ equal the range of the two outcomes (i.e. the absolute value of the difference
     of the largest and the smallest outcomes). Determine the p.m.f. $g(y)$ of $Y$ for $y=0, 1, 2, 3, 4, 5$.
    \end{enumerate}
\end{exercise}

\begin{solution}
題目表示出$X$是六面骰骰兩次中最小點數的隨機變數。因為合計有36種可能狀況，其中
可以知道數字為1的情況有11種，以此類推。更具體一點為，考慮$X=x$，那麼我們可以將 
\[
P(X\geq x)-P(X\geq x+1)=P(X=x)
\]
用來計算，因此
\[
(7-x)^2-(6-x)^2=13-2x
\]
所以$X$的p.m.f.為
\[
\frac{13-2x}{36}
\]

而$Y$的p.m.f.則可以直接計數(暫時沒想到其他辦法)得到
\[P(Y=y)=\begin{cases}
\frac{1}{6},\ y=0,\\
\frac{6-y}{18},\ y=1,2,3,4,5.
\end{cases}\]
\qed
\end{solution}

%%%%%%%%%%%%%%%%%%%%%%%%%%%%%%%%%%%%%%%%%%%

\begin{exercise}[2.1-11]
    In a lot (collection) of $100$ light bulbs, there are $5$ bad bulbs. An inspector
    inspects $10$ bulbs selected at random. Find the probability of finding at least 
    one defective bulb. Hint: First compute the probability of finding no defectives in 
    the sample.
\end{exercise}
\begin{solution}
    抽$10$個至少壞$1$個的機率
    $$
    P=1-\text{全部都是好的機率}=1-P'
    $$
    全部是好的機率就很好算了。
    $$
    P'=\frac{95}{100}\times\frac{94}{99}\cdots\times\frac{86}{91}
    $$
    藉由計算機，
    $$
    P=0.4163
    $$
    \qed
\end{solution}
%%%%%%%%%%%%%%%%%%%%%%%%%%%%%%%%%%%%%%%%%%%

\begin{exercise}[2.1-11]
    Let X be the no. of accidents per week in a factory. Let the p.m.f.
    of $X$ be
    $$
    f(x)=\frac{1}{(x+1)(x+2)},\quad x=0,1,2,\cdots.
    $$
    Find the conditional probability of $X\geq4$, given that $X\geq1$.
    Hint:Write $f(x)=1/(x+1)-1/(x+2)$.
\end{exercise}
\begin{solution}
    
    \qed
\end{solution}
%%%%%%%%%%%%%%%%%%%%%%%%%%%%%%%%%%%%%%%%%%%
\newpage
\bibliographystyle{siam}
\bibliography{ref.bib}
\nocite{*}


\end{document}
